\section{Evaluation Summary}\label{sec:evaluation-summary}

The evaluation supports five bounded conclusions.
First, the compiler expands compact source rules into much larger MIDI timelines, up to the 20{,}000-event cap from a
handful of source lines, confirming the structural-compression premise of the language.
Second, compile time decomposes cleanly across the pipeline: parsing and semantic analysis are negligible (under
$0.13$~ms on every fixture), and the total is governed by lowering and MIDI serialization.
Third, that total scales with the size of the generated score, lowering roughly linearly and MIDI serialization
slightly faster, consistent with the $O(E \log E)$ backend sort, so typical programs at the cap compile in about
$5$--$8$~ms.
Fourth, at a fixed output size cost depends on program shape through the lowering stage: per-note expression evaluation
is the worst case ($\sim$9~ms at the cap, roughly $2.7\times$ a flat loop), nesting depth is second, and chord density
and parallelism stay near the baseline.
Fifth, real compositions never approach these worst cases: the four measured pieces compile in $2.3$--$2.5$~ms,
dominated by process startup, with only $0.15$--$0.35$~ms of actual compiler work and modest $1.9$--$5.1$ events per
source line.
A separate comparison confirms that the version~2.0 type system adds no measurable compile-time cost over version~1.0
while producing byte-identical output.

These conclusions are bounded by the scope stated in Section~\ref{sec:evaluation-methodology} and
Section~\ref{sec:limits}; every figure is reproducible from a clean checkout via the \texttt{experiments/} suite.

The evidence gathered here closes the loop opened in the introduction.
The final chapter consolidates what the project contributes, states plainly what it does not claim, revisits the use
cases from Chapter~\ref{ch:introduction} against the delivered system, and sketches the extensions this work makes
possible.
